\section*{Asesmen}
\addcontentsline{toc}{section}{Asesmen}

Asesmen Bab I bersifat formatif dan tidak memiliki bobot tersendiri.
Nilainya menjadi bagian dari komponen \textbf{Kuis} (5\%) dan
\textbf{Tugas Teori} (8\%) pada RPS.

\begin{table}[H]
\centering
\caption{Instrumen asesmen dan indikator}
\label{tab:ch01-asesmen}
\begin{tabularx}{\textwidth}{@{}lXl@{}}
\toprule
\textbf{Instrumen} & \textbf{Indikator} & \textbf{Bobot} \\
\midrule
Kuis awal & Menjelaskan struktur buku, CPMK, dan alur belajar &
5\% (bagian) \\
Tugas rencana belajar & Menyusun rencana studi yang realistis dan
terukur & 8\% (bagian) \\
Praktikum 0 & Menyiapkan lingkungan dan menjalankan program contoh &
kelulusan prasyarat \\
\bottomrule
\end{tabularx}
\end{table}

Kriteria kelulusan Bab I: mahasiswa dapat menjelaskan peta buku,
menghubungkan bab dengan CPMK, dan menunjukkan bahwa lingkungan kerja
siap dipakai.
